\section{Conclusion}
\par This thesis focuses on the patch-based image focus detection task
\cite{Hou_2016_CVPR} and builds a unified experimental framework under
low-sample conditions, integrating image preprocessing, patch division,
automatic labeling, model training, and result evaluation into a complete
pipeline. On this basis, this thesis systematically compares the performance
of decision tree, random forest, support vector machine based on
Nystr\"{o}m kernel approximation, and a lightweight CNN
\cite{lecun2002gradient,krizhevsky2012imagenet} under the same data condition
and the same task definition, so as to analyze the applicability of different
modeling methods in this kind of task.

\par The core conclusion of this thesis is not to deny the advantage of CNNs in
image tasks, but to show that model selection still needs to be judged
according to specific task conditions. The experimental results show that the
lightweight CNN can obtain relatively strong overall recognition performance in
the binary task, but its training cost is also higher than that of classical
machine learning methods. At the same time, random forest and support vector
machine also show strong classification ability in multiple groups of
experiments, while requiring lower training cost and more direct
implementation. Especially in the low-sample, patch-level image focus
detection scenario studied in this thesis, both classical machine learning
methods and CNN are trained on a unified grayscale patch representation. The
experimental results show that under this input representation and task
definition, classical machine learning methods do not lose effectiveness simply
because they lack deep hierarchical representation. On the contrary, they show
relatively high practical value between effectiveness and efficiency. At the
same time, it should also be recognized that the training setting of the CNN
in this thesis itself contains a realistic trade-off between training time and
recognition effect. The main text does not take ``waiting for full
convergence'' as the only goal. Even under this premise, the CNN still shows
high accuracy and strong overall sharp-class recognition ability, but its time
cost remains clearly higher than that of classical machine learning methods.

\par Furthermore, the sample-size analysis shows that as the training ratio
increases from $20\%$ to $100\%$, the accuracy, sharp-class F1-score, and
macro F1 of most models improve only to a limited extent, while training time
continues to rise. This indicates that under the current task setting,
additional computational resources do not automatically turn into performance
gain of the same magnitude. It should be emphasized that this conclusion is
drawn for the scan over patch training sample ratios constructed from the fixed
set of original images in this thesis, and should not be directly generalized
to mean that increasing independent training data is always of limited benefit.
Within this scope, the results of this thesis further support the following
judgment: for image tasks, although CNNs have strong modeling ability, they are
not necessarily always the first choice. When task structure, data scale,
class-boundary characteristics, and resource constraints act together,
classical machine learning methods may still be more appropriate engineering
choices.

\par Overall, this thesis provides a reproducible comparison framework for
low-sample image focus detection and also offers a more targeted basis for
model selection in practical applications. For similar tasks, method selection
should consider recognition effect, training cost, data scale, and application
requirements together, rather than directly adopting deep learning solutions
only according to common experience from general vision tasks. It should also
be added that the supplementary experiments in this thesis already show that
the performance of the three-class CNN is clearly affected by training
strategies such as weighted sampling, and therefore the final evaluation of
different models should consider both model structure and training
configuration. Future work can further verify the conclusions of this thesis on
larger-scale data, higher-quality annotation, and more complex scenarios, so
as to examine whether increasing model complexity and computational resources
can bring more obvious practical gain when task difficulty continues to
increase.
